\section{Controlled Computational Characterization}
\label{sec:computational}

\subsection{Experimental question}

The computational study asks a narrower question than a conventional leaderboard: how does a mechanism-diverse set of optimizers behave across a controlled continuous black-box test space when problem identity, dimensionality, instance set, seed policy, and objective-evaluation budget are fixed in advance? The retained panel comprises Random Search, SciPy differential evolution (DE), bounded SciPy Nelder--Mead (NM), and pycma CMA-ES. These methods were chosen because their implementations and frozen configurations passed the publication acceptance path; the panel is intentionally smaller than the census.

All 24 noiseless BBOB functions are evaluated at dimensions \(D\in\{5,10,20,40\}\). Discovery uses instances 1--5 and held-out characterization uses instances 6--10. Stochastic configurations use seeds 11, 23, and 37, with checkpoints at \(100D\), \(300D\), and \(1000D\) objective evaluations. ObjectiveRecorder counts were cross-checked against COCO evaluation counts.

\subsection{Discovery and held-out COCO evidence}

Figure~\ref{fig:coco_discovery_heldout} places the official COCO ECDF views for discovery and held-out runs side by side for the 20-D all-function aggregate. The two panels are generated directly by \texttt{cocopp} from EF004 and EF005 artifacts, respectively; they are not reconstructed manually.

\begin{figure*}[H]
\centering
\begin{subfigure}[t]{0.49\textwidth}
\centering
\includegraphics[width=\linewidth]{figures/ef004_20d_all.pdf}
\caption{Discovery instances 1--5 (EF004).}
\label{fig:coco_discovery}
\end{subfigure}
\hfill
\begin{subfigure}[t]{0.49\textwidth}
\centering
\includegraphics[width=\linewidth]{figures/ef005_20d_all.pdf}
\caption{Held-out instances 6--10 (EF005).}
\label{fig:coco_heldout}
\end{subfigure}
\caption{Official COCO target-runtime ECDF views for the same four frozen configurations on all 24 noiseless BBOB functions at \(D=20\). The discovery and held-out panels are produced from disjoint instance sets under the same scientific configuration. The figure is used for condition-level characterization, not for a universal ranking or a post hoc binary failure boundary.}
\label{fig:coco_discovery_heldout}
\end{figure*}

The principal empirical observation is not a single global winner. The relative separation among methods varies with target difficulty and problem class, and the held-out panel preserves the same broad heterogeneity seen during discovery rather than collapsing to a stable total ordering. This is precisely why the paper reports ECDF/target-runtime behavior together with condition metadata instead of reducing the campaign to one average rank.

\subsection{Validated scale of the campaign}

Table~\ref{tab:campaign_scale} summarizes the promoted computational products. EF003 supplies fixed-budget cross-dimensional scaling; EF004 supplies algorithm-separated COCO target-runtime logging on discovery instances; EF005 repeats the target-runtime protocol on disjoint held-out instances.

\begin{table}[H]
\centering
\caption{Validated computational products used in the manuscript.}
\label{tab:campaign_scale}
\footnotesize
\begin{tabular}{p{0.11\linewidth}p{0.21\linewidth}p{0.18\linewidth}p{0.35\linewidth}}
\toprule
Freeze & Instance role & Scale & Publication use \\
\midrule
EF003 & Discovery, 1--5 & 1,152 dimension--algorithm--function--budget cells & Descriptive fixed-budget scaling at 100D, 300D, and 1000D \\
EF004 & Discovery, 1--5 & 14,400 checkpoint rows; 64 algorithm-labelled COCO roots & Official target-runtime/ECDF characterization \\
EF005 & Held-out, 6--10 & 14,400 checkpoint rows; 64 algorithm-labelled COCO roots & Held-out target-runtime/ECDF characterization \\
\bottomrule
\end{tabular}
\end{table}

Across EF004 and EF005, the combined validation logic required exact agreement between the final ObjectiveRecorder and COCO evaluation counts. Joint \texttt{cocopp} postprocessing completed for both campaigns. These checks do not establish algorithmic superiority; they establish that the performance views in Figure~\ref{fig:coco_discovery_heldout} are tied to validated evaluation accounting and algorithm-separated COCO roots.

\subsection{Fixed-budget and target-runtime views}

The study retains two complementary views. Fixed-budget analysis asks what has been achieved after a declared number of evaluations. Target-runtime analysis asks how many evaluations are required to reach a declared target, with unreached targets retained as unsuccessful observations. A method can be comparatively strong under one view and weaker under another without contradiction.

For target-runtime analysis, COCO's logger and postprocessing define the authoritative attainment semantics. Exact target runtimes are not reconstructed from sparse checkpoints, and private optimum values are not inferred from undocumented internals. For fixed-budget scaling, FE/D remains explicit so that the same nominal budget is not confused with equal wall-clock cost.

\subsection{Interpretation boundary}

The held-out campaign was designed to prevent the discovery instances from serving as their own sole confirmation. However, the original failure-frontier protocol required an exact numerical reliability threshold and target subset to be frozen before the held-out results were inspected. Those numerical choices were not fully frozen before EF005 execution. We therefore do not retrofit a binary frontier after observing the confirmation data.

This restriction is scientifically consequential. The experiments support condition-level target-runtime and scaling characterization across the evaluated BBOB conditions; they do not support a universal failure boundary, a universal optimizer ranking, or a mechanism-repair claim derived solely from deterioration. R12 was therefore closed as \texttt{NOT\_WARRANTED}, and no new optimizer is introduced.

\subsection{Cost and fairness}

Equal objective-evaluation budgets improve comparability but do not equalize computational cost. CMA-ES maintains covariance-related state, DE operates on a population, NM maintains simplex geometry, and Random Search has little optimizer-side state. Objective evaluations and optimizer overhead are therefore treated as different resources.

The comparator panel is similarly bounded. PSO and GA were not added merely to enlarge the experiment because their publication implementation/configuration provenance had not passed the same acceptance path when the panel was frozen. The paper favors fewer fully specified comparators over a broader but weaker comparison.

\subsection{Result synthesis}

The computational evidence supports three bounded conclusions. First, the same scientific configuration can be executed on disjoint discovery and held-out BBOB instances with exact objective-call accounting. Second, official COCO target-runtime products reveal substantial condition dependence, so a single aggregate rank would discard information central to interpretation. Third, the held-out campaign can validate a characterization claim while still being insufficient for a stronger threshold-dependent frontier claim. The empirical contribution is therefore not a new optimizer; it is an auditable demonstration of how claim strength changes when the experimental protocol is held fixed and the interpretation rule is enforced.
